\section{Discussion}
\label{sec:discussion}

\subsection{Which Decisions Actually Matter}
All recipes sit within $1.0$ macro-F1 point and, as
Sect.~\ref{sec:significance} shows, these gaps are not statistically
distinguishable under our tests (not proof of equivalence), so the ablation gives
only the direction of the effects.
Class-weighted focal loss on a plain head (M6) reached $0.6604$, matching the full
recipe within noise while giving the best minority recall and macro ROC-AUC (its
macro AP is second to M1's), so
the loss change gave the largest observed shift; the multi-view
head alone adds only $+0.30$ macro-F1.  More stages did not help: both
multi-stage systems were frozen
at a stage-1 checkpoint, so the extra stages never improved validation selection.
The threshold search collapsed to
$0.5$ (M3, M5, M6) or failed to transfer (M4), so the decision rule added no gain.

The accuracy leader is also the macro-F1 laggard (Sect.~\ref{sec:trap}), a direct
consequence of an 83\%-majority test set: each extra \textsc{clean} prediction
raises accuracy even when it misclassifies a minority comment, so an
accuracy-optimised system would under-report offensive and hateful content---the
opposite of the product requirement.  Selection should therefore use macro-F1 (or
minority recall), driven by validation so the test split stays honest.

\subsection{Threats to Validity}
\label{sec:threats}
Each system was trained once (seed 42), so the estimates carry run-to-run variance
the shared split and seed do not remove; the bootstrap intervals quantify sampling
noise but not training noise, and as the observed gaps already lie inside one
interval we claim no recipe is reliably better.  The official split reuses comments
across train/dev/test (Sect.~\ref{sec:dataset}); we report this rather than hide it,
and the de-duplicated check leaves the macro-F1 ordering effectively unchanged
(Table~\ref{tab:duplicate-free}), but absolute scores are
optimistic under a strictly leakage-free protocol.  ViHSD is the only corpus, its
external baselines (Table~\ref{tab:external}) reprint published numbers under other
settings, and the latency runs were inconsistent between identical architectures,
so no cost claim is made.  Agreement is only moderate ($\kappa\approx0.5$), so part
of the error is irreducible.

\subsection{Implications}
For applied moderation the message is negative but worth reporting: no single trick
solves Vietnamese hate speech detection, and added complexity---curricula, searched
thresholds, longer training---did not pay off.  The
minority classes remain hard
(OFFENSIVE F1 stuck near $0.46$ across all six systems), which is consistent with
the labelled data, not the optimiser, being the binding constraint; more labelled
offensive examples or span-level supervision as in ViHOS \cite{hoang2023vihos}
are plausible next steps.
